\setcounter{chapter}{7}
\setcounter{section}{0}
\setcounter{table}{0}
\setcounter{figure}{0}

\addcontentsline{toc}{chapter}{Apéndice 2: Procedimiento reproducible del contenedor y CUDA}

\makeatletter
\def\@currentlabel{\thechapter}
\makeatother
\label{ape:contenedor-cuda}

\chapter*{Apéndice 2: Procedimiento reproducible del contenedor y CUDA}

Este apéndice resume la creación del contenedor LXC con acceso a las GPUs y la instalación del stack CUDA 6.5, asumiendo que el driver 340.108 ya está operativo en el host (apéndice~\ref{ape:host}).

\section{Creación y configuración del contenedor}

\begin{Verbatim}
# 1. Obtener la plantilla (consultar el nombre exacto con pveam available)
pveam update
pveam available --section system | grep debian-12
pveam download local <nombre-exacto-de-la-plantilla>

# 2. Crear el contenedor (privilegiado, con nesting)
pct create 200 local:vztmpl/<plantilla>.tar.zst \
  --hostname gpu-tfg \
  --cores 4 \
  --memory 4096 \
  --rootfs local-lvm:32 \
  --net0 name=eth0,bridge=vmbr0,ip=dhcp \
  --unprivileged 0 \
  --features nesting=1

# 3. Configurar el paso de dispositivos en /etc/pve/lxc/200.conf
#    lxc.cgroup2.devices.allow: c 195:* rwm
#    lxc.mount.entry: /dev/nvidia0 dev/nvidia0 none bind,optional,create=file
#    lxc.mount.entry: /dev/nvidia1 dev/nvidia1 none bind,optional,create=file
#    lxc.mount.entry: /dev/nvidiactl dev/nvidiactl none bind,optional,create=file

# 4. Verificar los nodos en el host y arrancar
nvidia-smi >/dev/null
ls -la /dev/nvidia0 /dev/nvidia1 /dev/nvidiactl
pct start 200
pct enter 200
\end{Verbatim}

\section{Stack userspace NVIDIA en el contenedor}

\begin{Verbatim}
apt update
apt install -y wget build-essential

cd ~
wget https://us.download.nvidia.com/XFree86/Linux-x86_64/340.108/NVIDIA-Linux-x86_64-340.108.run
chmod +x NVIDIA-Linux-x86_64-340.108.run
./NVIDIA-Linux-x86_64-340.108.run --no-kernel-module --no-x-check
#   -> declinar librerías de 32 bits y configuración X11
nvidia-smi   # ambas GPUs visibles
\end{Verbatim}

\section{CUDA Toolkit 6.5}

\begin{Verbatim}
# 1. Descarga del meta-instalador
cd ~
wget https://developer.download.nvidia.com/compute/cuda/6_5/rel/installers/cuda_6.5.14_linux_64.run
chmod +x cuda_6.5.14_linux_64.run

# 2. Nivel 1: extraer los subinstaladores
mkdir -p ~/cuda-extract
./cuda_6.5.14_linux_64.run --extract=$HOME/cuda-extract

# 3. Nivel 2: extraer el toolkit a ruta persistente y ejecutar el instalador Perl
cd ~/cuda-extract
chmod +x cuda-linux64-rel-6.5.14-18749181.run
./cuda-linux64-rel-6.5.14-18749181.run --keep --target ~/cuda-toolkit-extracted --noexec
cd ~/cuda-toolkit-extracted
perl -I. ./install-linux.pl -noprompt -prefix=/usr/local/cuda-6.5

# 4. Samples (extracción manual y copia del árbol)
cd ~/cuda-extract
chmod +x cuda-samples-linux-6.5.14-18745345.run
./cuda-samples-linux-6.5.14-18745345.run --keep --target ~/cuda-samples-extracted --noexec
mkdir -p ~/NVIDIA_CUDA-6.5_Samples
cp -r ~/cuda-samples-extracted/cuda-samples/* ~/NVIDIA_CUDA-6.5_Samples/

# 5. Configuración del entorno
echo 'export PATH=/usr/local/cuda-6.5/bin:$PATH' >> ~/.bashrc
echo 'export LD_LIBRARY_PATH=/usr/local/cuda-6.5/lib64:$LD_LIBRARY_PATH' >> ~/.bashrc
source ~/.bashrc
nvcc --version
\end{Verbatim}

\section{Validación y benchmarks}

\begin{Verbatim}
# 1. deviceQuery (cómputo real)
find ~/NVIDIA_CUDA-6.5_Samples -iname "deviceQuery" -type d
cd <ruta>/1_Utilities/deviceQuery
make && ./deviceQuery

# 2. Compilación de los benchmarks en el chroot Debian Jessie (solo compilación)
apt install -y debootstrap
debootstrap --no-check-gpg --arch=amd64 jessie /root/old-chroot http://archive.debian.org/debian/
#  (sources.list con archive.debian.org, build-essential, bind mounts de /proc,
#   /dev, /sys, el toolkit y los samples; ver capítulo 5)

# 3. Ejecución de los benchmarks en el contenedor principal
cd ~/NVIDIA_CUDA-6.5_Samples/0_Simple/matrixMul
./matrixMul --device=0
./matrixMul --device=1
cd ~/NVIDIA_CUDA-6.5_Samples/1_Utilities/bandwidthTest
./bandwidthTest --device=0
./bandwidthTest --device=1
\end{Verbatim}
